\paragraph{window-\texorpdfstring{$6$}{6} 的内禀自校准与边界证书：末位均值比率、四块谱数域、返时母函数与半群阈值}
本段在不引入外部标定的前提下，把若干与 window-\(6\) 锚点直接相连、且可由有限数据审计的结论压缩为一条可独立引用的链条：其一是 bin-fold 两态一致渐近所诱导的 \(\varphi\) 自校准统计量；其二是 boundary uplift 的四步平移同构所给出的家族数窗口量子化与最小窗唯一性；其三是四块粗粒化核的不可约三次谱数域与 \(S_3\) 伽罗瓦群；其四是边界块 \(U_R\) 的首返时间分布闭式与通量倒数律；其五是由 tail-offset 三元组 \(\{21,34,55\}\) 所生成数值半群的 Frobenius 阈值与由此导出的“维数筛法有限性”。
\par

\begin{theorem}[黄金比的内禀自校准：由末位条件均值比率恢复 \texorpdfstring{$\varphi$}{phi}]\label{thm:xi-foldbin-intrinsic-self-calibration-phi-lastbit-means}
\leanverified{paper\_xi\_foldbin\_intrinsic\_self\_calibration\_phi\_lastbit\_means}
在 bin-fold 两态一致渐近（定理 \ref{thm:fold-bin-two-state-asymptotic}）成立的口径下，
令 \(d_m^{\mathrm{bin}}(x):=\bigl|(\mathrm{Fold}^{\mathrm{bin}}_m)^{-1}(x)\bigr|\)（\(x\in X_m\)），并按末位划分
\[
X_m^{(0)}:=\{x\in X_m:\ x_m=0\},
\qquad
X_m^{(1)}:=\{x\in X_m:\ x_m=1\}.
\]
定义可观测统计量
\[
\widehat\varphi_m
\ :=\
\frac{\frac{1}{|X_m^{(0)}|}\sum_{x\in X_m^{(0)}} d_m^{\mathrm{bin}}(x)}{\frac{1}{|X_m^{(1)}|}\sum_{x\in X_m^{(1)}} d_m^{\mathrm{bin}}(x)}.
\]
则存在与 \(m\) 无关常数 \(C>0\) 使得
\[
\bigl|\widehat\varphi_m-\varphi\bigr|\ \le\ C\left(\frac{\varphi}{2}\right)^m,
\]
从而 \(\widehat\varphi_m\) 以指数精度收敛到 \(\varphi\)。
该结论给出一条无需外部标定的自校准机制：仅凭折叠纤维多重度与末位可见位即可恢复增长常数。
\end{theorem}

\begin{proof}
由定理 \ref{thm:fold-bin-two-state-asymptotic}，存在常数 \(C_0>0\) 使对充分大 \(m\) 与任意 \(x\in X_m\) 有一致估计
\[
\Bigl|d_m^{\mathrm{bin}}(x)-\varphi^{-x_m}\Bigl(\frac{2}{\varphi}\Bigr)^m\Bigr|\le C_0.
\]
对 \(X_m^{(0)}\)、\(X_m^{(1)}\) 分别取均值，得
\[
\frac{1}{|X_m^{(0)}|}\sum_{x\in X_m^{(0)}} d_m^{\mathrm{bin}}(x)
=\Bigl(\frac{2}{\varphi}\Bigr)^m+O(1),
\qquad
\frac{1}{|X_m^{(1)}|}\sum_{x\in X_m^{(1)}} d_m^{\mathrm{bin}}(x)
=\varphi^{-1}\Bigl(\frac{2}{\varphi}\Bigr)^m+O(1),
\]
其中 \(O(1)\) 的常数与 \(m\) 无关。
两式相除并用 \((2/\varphi)^{-m}=(\varphi/2)^m\) 得
\[
\widehat\varphi_m=\varphi\cdot\frac{1+O((\varphi/2)^m)}{1+O((\varphi/2)^m)}
=\varphi+O\!\left(\left(\frac{\varphi}{2}\right)^m\right),
\]
从而存在常数 \(C\) 满足所示误差界。证毕。
\end{proof}

\begin{theorem}[家族数的窗口量子化与最小窗唯一性：\texorpdfstring{$N_f=|X_m^{\mathrm{bdry}}|$}{Nf=|Xbdry|}]\label{thm:xi-family-number-window-quantization-minimal-window}
\leanverified{paper\_xi\_family\_number\_window\_quantization\_minimal\_window}
固定每代 \(SU(5)\) 手征计数 \(15\) 为外部输入，并要求 “\(SO(10)\)-型最小补全” 在同一窗口 \(m\) 上由边界 singlet 通道完成，即满足离散恒等式
\[
15N_f+|X_m^{\mathrm{bdry}}|=16N_f.
\]
则必有
\[
N_f=|X_m^{\mathrm{bdry}}|=F_{m-2},
\]
从而家族数在该制度下只能取 Fibonacci 值。
特别地，若要求 \(N_f=3\)，则唯一解为 \(m=6\)（最小窗唯一性），并与 window-\(6\) uplift 三分支中 \(F_9=34\) 分支的离散对齐一致（表 \ref{tab:window6_family_uplift_lock} 与 \eqref{eq:window6_family_singlet_completion}）。
\end{theorem}

\begin{proof}
由所给离散恒等式立得 \(N_f=|X_m^{\mathrm{bdry}}|\)。
另一方面，由定理 \ref{thm:boundary-shift4-uplift-isomorphism} 的 uplift 同构 \(X_{m-4}\simeq X_m^{\mathrm{bdry}}\)（\(m\ge 6\)）与 \(|X_n|=F_{n+2}\) 得
\(|X_m^{\mathrm{bdry}}|=|X_{m-4}|=F_{m-2}\)，从而得到 \(N_f=F_{m-2}\)。
若 \(N_f=3\)，则 \(F_{m-2}=3\) 强制 \(m-2=4\)，即 \(m=6\)，且该解唯一。证毕。
\end{proof}

\begin{theorem}[window-\(6\) 四块粗粒化核的谱数域刚性：不可约三次域与 \texorpdfstring{$S_3$}{S3} 伽罗瓦群]\label{thm:xi-window6-four-block-kernel-cubic-field-s3}
\leanverified{paper\_xi\_window6\_four\_block\_kernel\_cubic\_field\_s3}
令四块骨架的粗粒化 Markov 核 \(T\) 如 \eqref{eq:window6_edge_flux_skeleton}，
并记其特征多项式分解为
\[
\chi_T(t)=(t-1)\cdot p(t),
\qquad
p(t)=\frac{48114t^3+7263t^2-506t-55}{48114}.
\]
则：
\begin{enumerate}
  \item 三次因子 \(p(t)\) 在 \(\mathbb Q[t]\) 上不可约；
  \item 其判别式 \(\Delta=\mathrm{Disc}(48114t^3+7263t^2-506t-55)\) 非平方（见 \eqref{eq:window6_edge_flux_skeleton_arithmetic}）；
  \item 因而分裂域的伽罗瓦群为 \(S_3\)，并且 \(T\) 的全部非平凡本征值均为代数次数 \(3\) 的实数。
\end{enumerate}
因此非平凡谱生成一个不可交换的三次数域 \(K_T=\mathbb Q(\lambda)\)，并且“粗粒化混合”所需的最小代数复杂度在数域意义下为三次且不可降阶。
\end{theorem}

\begin{proof}
对三次多项式，\(\mathbb Q[t]\) 上不可约当且仅当无有理根。
把整系数多项式
\(
f(t):=48114t^3+7263t^2-506t-55
\)
模 \(13\) 化简，得
\(
f(t)\equiv t^3+9t^2+t+10\pmod{13}
\)。
直接检验 \(t\in\mathbb F_{13}\) 的全部取值可知该三次多项式在 \(\mathbb F_{13}\) 上无根，故在 \(\mathbb F_{13}[t]\) 上不可约，从而 \(f\) 在 \(\mathbb Q[t]\) 上不可约，\(p\) 亦不可约。
\par
由 \eqref{eq:window6_edge_flux_skeleton_arithmetic}，判别式 \(\Delta\) 非平方。
对不可约三次多项式，其伽罗瓦群为 \(A_3\) 当且仅当判别式为平方；否则为 \(S_3\)。
因此本例伽罗瓦群为 \(S_3\)。
又因 \(\Delta>0\)，三根互异且均为实数，从而 \(T\) 的三个非平凡本征值均为代数次数 \(3\) 的实数。证毕。
\end{proof}

\begin{corollary}[判别式分解的审计素数：\texorpdfstring{$23$}{23} 的不可消去出现]\label{cor:xi-window6-four-block-kernel-disc-audit-prime-23}
沿定理 \ref{thm:xi-window6-four-block-kernel-cubic-field-s3} 的整系数多项式
\(
f(t)=48114t^3+7263t^2-506t-55
\)。
其判别式满足
\[
\boxed{\;
\Disc(f)
=2^{2}\cdot 3^{2}\cdot 11^{2}\cdot 13^{2}\cdot 23\;}
\]
并且 \(\det(tI-T)\) 的常数项分解含因子
\(
23\cdot 11\cdot 5
\)
（见证书 \eqref{eq:window6_edge_flux_skeleton_arithmetic}）。
因此素数 \(23\) 作为四块粗粒化核的代数数论指纹在该口径下不可被规避或坐标变换消去：其出现由判别式与常数项同时锁定。
\end{corollary}

\begin{corollary}[粗粒化谱数域与黄金尺度域的互素分离：合成域次数为 \(6\)]\label{cor:xi-window6-field-disjointness-degree-six}
令 \(K_T=\mathbb Q(\lambda)\) 为定理 \ref{thm:xi-window6-four-block-kernel-cubic-field-s3} 的非平凡谱所生成三次数域，
并令 \(K_\varphi=\mathbb Q(\varphi)\) 为黄金二次数域。
则必有
\[
K_T\cap K_\varphi=\mathbb Q,
\]
从而合成域次数满足
\[
[K_TK_\varphi:\mathbb Q]=[K_T:\mathbb Q]\,[K_\varphi:\mathbb Q]=6.
\]
\end{corollary}

\begin{proof}
交域 \(K_T\cap K_\varphi\) 是 \(K_T\) 的子域，其次数必须整除 \([K_T:\mathbb Q]=3\)。
另一方面，\(K_T\cap K_\varphi\subseteq K_\varphi\) 且 \([K_\varphi:\mathbb Q]=2\)，故其次数亦整除 \(2\)。
于是 \([K_T\cap K_\varphi:\mathbb Q]\) 同时整除 \(2\) 与 \(3\)，只能为 \(1\)，即交域为 \(\mathbb Q\)。
线性无关性推出合成域次数为积。证毕。
\end{proof}

\begin{corollary}[混合指数与模尺度的非共振：\texorpdfstring{$-\log\Lambda/\log\varphi\notin\mathbb Q$}{-logLambda/logphi notin Q}]\label{cor:xi-window6-logLambda-logphi-irrational}
令
\[
\Lambda:=\max\{|\lambda|:\ \lambda\in\mathrm{spec}(T)\setminus\{1\}\}\in(0,1)
\]
为四块粗粒化核 \(T\) 的非平凡谱半径（取绝对值的谱半径）。
则
\[
\frac{-\log\Lambda}{\log\varphi}\notin\mathbb Q.
\]
\end{corollary}

\begin{proof}
由定理 \ref{thm:xi-window6-four-block-kernel-cubic-field-s3}，\(T\) 的非平凡本征值为代数次数 \(3\) 的实数；
因此 \(\Lambda\) 等于某个非平凡实本征值 \(\lambda_\star\) 的绝对值，故 \(\Lambda\in K_T\setminus\mathbb Q\)。
若 \(-\log\Lambda/\log\varphi=p/q\in\mathbb Q\)（既约，\(q\ge 1\)），则
\(
\Lambda^q=\varphi^{-p}\in K_\varphi
\)。
但 \(\Lambda^q\in K_T\)，故 \(\Lambda^q\in K_T\cap K_\varphi=\mathbb Q\)（推论 \ref{cor:xi-window6-field-disjointness-degree-six}）。
由于 \(\Lambda\in(0,1)\) 得 \(p\neq 0\)，而 \(\varphi^{-p}\notin\mathbb Q\)，矛盾。证毕。
\end{proof}

\begin{proposition}[Baker 型下界在本体系中的专门化：尺度拍频的有效 Diophantine 约束]\label{prop:xi-window6-baker-effective-beat-bound}
在推论 \ref{cor:xi-window6-logLambda-logphi-irrational} 的设定下，令 \(\alpha_1=\Lambda\)、\(\alpha_2=\varphi\)。
则由对数线性形式下界理论（Baker--W\"ustholz/Matveev 型结果）推出：存在可有效取值的常数 \(C>0\)，使对任意整数对 \((n,k)\neq(0,0)\) 有
\[
\bigl|n\log\Lambda+k\log\varphi\bigr|
\ \ge\
\exp\!\bigl(-C\log(1+\max\{|n|,|k|\})\bigr).
\]
因此任意“混合尺度 \(\Lambda^n\)”与“模尺度 \(\varphi^{-k}\)”之间的近共振误差至多以多对数速度出现，排除了指数级同步。
\end{proposition}

\begin{proof}
由推论 \ref{cor:xi-window6-logLambda-logphi-irrational}，\(\Lambda\) 与 \(\varphi\) 乘法独立。
对数线性形式下界给出对任意非零整数线性组合 \(n\log\Lambda+k\log\varphi\) 的有效下界，其量级为 \((1+\max\{|n|,|k|\})^{-C}\)（等价写作所示指数形式）；详见 \cite{BakerWustholz2007LogarithmicForms,Matveev2000LinearFormsLogs}。证毕。
\end{proof}

\paragraph{四块商图的 Kirchhoff 复杂度与粗粒化 Green 核的分母障碍}
在四块骨架证书 \eqref{eq:window6_edge_flux_skeleton} 的记号下，令 \(\mathcal{G}_{\mathrm{blk}}\) 为加权无向多重图：顶点为 \(\{U_2,U_1,U_L,U_R\}\)，且当 \(\alpha\neq\beta\) 时边重数为 \(E_{\alpha\beta}\)（忽略块内自环）。
记其加权 Laplace 矩阵为 \(L_{\mathrm{blk}}\)（对角为度数、非对角为负边重数），并令 \(L_{\mathrm{blk}}^{(\alpha)}\) 为删去同一行列所得主余子式。

\begin{theorem}[回路格的电阻型 Gram 行列式与生成树权]\label{thm:xi-flow-lattice-gram-determinant-tree-weight}
\leanverified{paper\_xi\_flow\_lattice\_gram\_determinant\_tree\_weight}
令 \(G=(V,E)\) 为有限连通（允许多重边）无向图，并固定一组定向以定义去掉一行的有向关联矩阵
\(
B\in\ZZ^{(|V|-1)\times|E|}
\)
（于是 \(\mathrm{rank}\,B=|V|-1\)）。
给定边权 \(w_e>0\)，令 \(W:=\mathrm{diag}(w_e)_{e\in E}\)。
定义整数流格
\[
\Lambda_{\mathrm{fl}}:=\ker(B)\cap\ZZ^{E}\subset \RR^E,
\]
并在 \(\RR^E\) 上取电阻型二次型
\[
\langle x,y\rangle_{W^{-1}}:=x^{\mathsf T}W^{-1}y.
\]
记 \(\tau_w(G)\) 为加权生成树总权（对每棵生成树取边权乘积并求和）。
则 \(\Lambda_{\mathrm{fl}}\) 在该二次型下的 Gram 行列式满足精确恒等式
\[
\det\!\bigl(\mathrm{Gram}_{W^{-1}}(\Lambda_{\mathrm{fl}})\bigr)=\frac{\tau_w(G)}{\det(W)}.
\]
等价地，其协体积为
\[
\mathrm{covol}_{W^{-1}}(\Lambda_{\mathrm{fl}})=\sqrt{\tau_w(G)/\det(W)}.
\]
\end{theorem}
\begin{proof}
取任一生成树 \(T\subseteq E\)，令弦边集 \(C:=E\setminus T\)，并按 \(E=T\sqcup C\) 重排列边。
相应写 \(B=[B_T\ B_C]\) 与 \(W=\mathrm{diag}(W_T,W_C)\)。
由于 \(T\) 为生成树，\(B_T\in\ZZ^{(|V|-1)\times(|V|-1)}\) 为可逆整矩阵且 \(\det(B_T)=\pm 1\)。
令
\[
C_*:=\begin{pmatrix}-B_T^{-1}B_C\\ I\end{pmatrix}\in\ZZ^{E\times|C|},
\]
则 \(BC_*=0\)，且其列向量给出 \(\Lambda_{\mathrm{fl}}\) 的一组整基。
故 Gram 矩阵为
\[
G_*:=C_*^{\mathsf T}W^{-1}C_*
=W_C^{-1}+B_C^{\mathsf T}B_T^{-\mathsf T}W_T^{-1}B_T^{-1}B_C.
\]
由矩阵行列式引理与 Sylvester 恒等式 \(\det(I+AB)=\det(I+BA)\)，得
\[
\det(G_*)
=\det(W_C^{-1})\cdot
\det\!\Bigl(I+(B_TW_TB_T^{\mathsf T})^{-1}B_CW_CB_C^{\mathsf T}\Bigr).
\]
另一方面，加权约化拉普拉斯为
\[
L_w:=BW B^{\mathsf T}=B_TW_TB_T^{\mathsf T}+B_CW_CB_C^{\mathsf T},
\]
并有分解
\[
\det(L_w)=\det(B_TW_TB_T^{\mathsf T})\cdot
\det\!\Bigl(I+(B_TW_TB_T^{\mathsf T})^{-1}B_CW_CB_C^{\mathsf T}\Bigr).
\]
由于 \(\det(B_TW_TB_T^{\mathsf T})=\det(W_T)\)（因 \(\det(B_T)=\pm1\)），比较两式得到
\[
\det(G_*)
=\det(W_C^{-1})\cdot\frac{\det(L_w)}{\det(W_T)}
=\frac{\det(L_w)}{\det(W_T)\det(W_C)}
=\frac{\det(L_w)}{\det(W)}.
\]
加权矩阵树定理给出 \(\det(L_w)=\tau_w(G)\)，从而结论成立。证毕。
\end{proof}

\begin{corollary}[固定边界偏差后的整数流扭量协体积]\label{cor:xi-affine-flow-torsor-covolume}
\leanverified{paper\_xi\_affine\_flow\_torsor\_covolume}
在定理 \ref{thm:xi-flow-lattice-gram-determinant-tree-weight} 的设定下，给定任意 \(b\in\ZZ^{|V|-1}\) 使得同余方程 \(BF=b\) 在 \(\ZZ^E\) 中可解。
取一解 \(F^{(0)}\in\ZZ^E\)，并令
\[
S(b):=\{F\in\ZZ^E:\ BF=b\}=F^{(0)}+\Lambda_{\mathrm{fl}}.
\]
则 \(S(b)\) 为以 \(\Lambda_{\mathrm{fl}}\) 为平移群的仿射格。
因此在二次型 \(\langle\cdot,\cdot\rangle_{W^{-1}}\) 下，解集的自由度密度由
\[
\mathrm{covol}_{W^{-1}}(\Lambda_{\mathrm{fl}})=\sqrt{\tau_w(G)/\det(W)}
\]
精确控制。
\end{corollary}
\begin{proof}
若 \(F\in S(b)\)，则 \(B(F-F^{(0)})=0\)，即 \(F-F^{(0)}\in\Lambda_{\mathrm{fl}}\)，反之亦然，故 \(S(b)=F^{(0)}+\Lambda_{\mathrm{fl}}\)。
协体积等式由定理 \ref{thm:xi-flow-lattice-gram-determinant-tree-weight} 直接给出。证毕。
\end{proof}

\begin{corollary}[window-\(6\) 四块商图的加权树数分解与审计素数 \texorpdfstring{$p_\star=571$}{p*=571}]\label{cor:xi-window6-block-tree-number-audit-prime}
\leanverified{paper\_xi\_window6\_block\_tree\_number\_audit\_prime}
对任意块 \(\alpha\in\{2,1,L,R\}\)，主余子式的行列式一致且等于同一加权树数
\[
\det\!\bigl(L_{\mathrm{blk}}^{(\alpha)}\bigr)
=\tau(\mathcal{G}_{\mathrm{blk}})
=123336
=2^{3}\cdot 3^{3}\cdot 571.
\]
因此在该四块几何的 Kirchhoff 复杂度中出现唯一素因子 \(p_\star:=571\) 满足 \(\gcd(p_\star,6)=1\)；该素数可作为块级树复杂度意义下的算术指纹。
\end{corollary}

\begin{proof}
矩阵树定理给出 \(\det(L_{\mathrm{blk}}^{(\alpha)})=\tau(\mathcal{G}_{\mathrm{blk}})\) 且与 \(\alpha\) 无关。
数值与分解由证书 \eqref{eq:window6_edge_flux_skeleton_arithmetic} 给出。证毕。
\end{proof}

\begin{proposition}[粗粒化 Green 核的算术闭包与 \texorpdfstring{$p_\star$}{p*} 的最小性]\label{prop:xi-window6-green-kernel-arithmetic-audit-prime}
在四块粗粒化核 \(T\)（\eqref{eq:window6_edge_flux_skeleton}）与平稳分布
\(
\pi_\alpha:=|U_\alpha|/64
\)
下，定义基本 Green 核
\[
Z:=(I-T+\mathbf 1\pi^\top)^{-1}.
\]
则 \(Z\) 的全部矩阵元属于 \(\mathbb{Z}[1/2,1/3,1/p_\star]\)，并且在既约分母中除 \(2,3\) 外不出现其他素因子。
此外，
\[
\det(I-T+\mathbf 1\pi^\top)=\frac{2^{4}\,p_\star}{3^{6}\cdot 11}
\]
（证书 \eqref{eq:window6_green_kernel_audit_prime_571}），从而 \(p_\star\) 在 \(Z\) 中不可被整体消去：存在矩阵元的既约分母含素因子 \(p_\star\)（见表 \ref{tab:window6_green_kernel_Z}）。
\end{proposition}

\begin{proof}
由脚本审计输出 \eqref{eq:window6_green_kernel_audit_prime_571} 与表 \ref{tab:window6_green_kernel_Z}，\(Z\) 的分母谱仅含 \(2,3,p_\star\)。
又由 Cramer 法则，\(Z\) 的任一矩阵元为 \(\mathrm{adj}(I-T+\mathbf 1\pi^\top)\) 的余子式除以 \(\det(I-T+\mathbf 1\pi^\top)\)。
由于该行列式既约分子含 \(p_\star\)，且表 \ref{tab:window6_green_kernel_Z} 中确有既约分母含 \(p_\star\) 的矩阵元，故 \(p_\star\) 不能在 Green 核层面被消去。证毕。
\end{proof}

\begin{theorem}[Markov--导数行列式恒等式：坏素数当且仅当特征多项式在 \(1\) 处模 \(p\) 重根]\label{thm:xi-markov-derivative-determinant-bad-prime}
\leanverified{paper\_xi\_markov\_derivative\_determinant\_bad\_prime}
设 \(T\) 为 \(N\times N\) 行随机矩阵，\(\pi^\top\) 为其平稳分布（\(\pi^\top T=\pi^\top\)，\(\pi^\top\mathbf 1=1\)），并假设特征值 \(1\) 在 \(\mathbb Q\) 上为单根。
定义
\[
A:=I-T+\mathbf 1\pi^\top,\qquad Z:=A^{-1},
\qquad
\chi_T(t):=\det(tI-T).
\]
则：
\begin{enumerate}
  \item 存在严格恒等式
  \[
  \det(A)=\chi_T'(1).
  \]
  \item 若 \(T,\pi\) 的分量均为有理数，取一正整数 \(D\) 使 \(DA\in M_N(\mathbb Z)\)，则对任意素数 \(p\nmid D\) 以下条件等价：
  \[
  DA\ \text{在 }\mathbb F_p\text{ 上奇异}
  \quad\Longleftrightarrow\quad
  \chi_T(t)\equiv (t-1)^2q(t)\pmod p.
  \]
  并且一旦成立，则 \(p\) 必出现在 \(Z\) 的既约分母中，即 \(Z\notin M_N(\mathbb Z[1/D])\)。
\end{enumerate}
\end{theorem}
\begin{proof}
令 \(V=\mathbb Q^N\)。
由 \(\pi^\top T=\pi^\top\) 可知 \(\ker(\pi^\top)\) 为 \(T\)-不变子空间；且 \(V=\mathbb Q\mathbf 1\oplus \ker(\pi^\top)\) 为直和分解。
对 \(v\in \ker(\pi^\top)\)，有 \(\mathbf 1\pi^\top v=0\)，故 \(Av=(I-T)v\)；
对 \(\mathbf 1\)，有 \(A\mathbf 1=\mathbf 1-T\mathbf 1+\mathbf 1\pi^\top\mathbf 1=\mathbf 1\)。
因此
\[
\det(A)=\det(I-T\mid_{\ker(\pi^\top)}).
\]
另一方面，由上述直和分解，特征多项式分解为
\[
\chi_T(t)=(t-1)\cdot \det(tI-T\mid_{\ker(\pi^\top)}),
\]
从而
\(
\chi_T'(1)=\det(I-T\mid_{\ker(\pi^\top)})=\det(A)
\)，得到 (1)。
\par
对 (2)，当 \(p\nmid D\) 时，\(DA\) 在 \(\mathbb F_p\) 上奇异当且仅当 \(\det(A)\equiv 0\pmod p\)。
由 (1) 这等价于 \(\chi_T'(1)\equiv 0\pmod p\)，又因 \(\chi_T(1)=0\)，故等价于 \(t=1\) 在 \(\chi_T(t)\) 中为模 \(p\) 的重根。
最后 \(Z=A^{-1}=\mathrm{adj}(A)/\det(A)\)，若 \(\det(A)\) 的既约分子被 \(p\) 整除，则 \(Z\) 的既约分母必含素因子 \(p\)。证毕。
\end{proof}

\begin{table}[H]
\centering
\scriptsize
\setlength{\tabcolsep}{6pt}
\caption{Window-6 four-block fundamental matrix $Z=(I-T+\mathbf 1\pi^\top)^{-1}$ (exact rationals).}
\label{tab:window6_green_kernel_Z}
\begin{tabular}{c c c c c}
\toprule
 & $U_2$ & $U_1$ & $U_L$ & $U_R$ \\
\midrule
$U_2$ & $(\frac{546339}{584704})$ & $(\frac{32153}{877056})$ & $(\frac{1441}{584704})$ & $(\frac{23233}{877056})$ \\
$U_1$ & $(\frac{26307}{584704})$ & $(\frac{288211}{292352})$ & $(\frac{9345}{584704})$ & $(-\frac{13685}{292352})$ \\
$U_L$ & $(\frac{4323}{584704})$ & $(\frac{34265}{877056})$ & $(\frac{548833}{584704})$ & $(\frac{13057}{877056})$ \\
$U_R$ & $(\frac{69699}{584704})$ & $(-\frac{150535}{877056})$ & $(\frac{13057}{584704})$ & $(\frac{903457}{877056})$ \\
\bottomrule
\end{tabular}
\end{table}


\begin{proposition}[块内停留指数四元组与 \texorpdfstring{$22/7$}{22/7} 的刚性出现]\label{prop:xi-window6-stay-indices-22-7}
\leanverified{paper\_xi\_window6\_stay\_indices\_22\_7}
令块内停留概率 \(s_\alpha:=T_{\alpha\alpha}\)，并定义停留指数 \(I_\alpha:=s_\alpha^{-1}\)。
则
\[
(I_2,I_1,I_L,I_R)=\left(\frac{81}{28},\ \frac{22}{7},\ \frac{27}{2},\ 9\right).
\]
特别地，\(I_1=22/7\) 以严格等式出现，且仅由离散计数 \((|U_1|,E_{11})=(22,21)\) 与定义 \(s_1=2E_{11}/(6|U_1|)\) 强制，不依赖任何连续调参。
\end{proposition}

\begin{proof}
由 \eqref{eq:window6_edge_flux_skeleton}，
\(
s_2=\frac{28}{81},\ s_1=\frac{7}{22},\ s_L=\frac{2}{27},\ s_R=\frac{1}{9}
\)。
逐项取倒数即得。证毕。
\end{proof}

\leanverified{paper\_xi\_window6\_kemeny\_constant\_rational}
\begin{proposition}[Kemeny 常数的有理闭式与 \texorpdfstring{$p'(1)/p(1)$}{p'(1)/p(1)} 恒等式]\label{prop:xi-window6-kemeny-constant-rational}
记特征多项式分解 \(\chi_T(t)=(t-1)\,p(t)\)，其中
\[
p(t)=\frac{48114t^{3}+7263t^{2}-506t-55}{48114}.
\]
则 Kemeny 常数
\(
K:=\sum_{\lambda\in\mathrm{Spec}(T)\setminus\{1\}}\frac{1}{1-\lambda}
\)
满足严格恒等式
\[
K=\frac{p'(1)}{p(1)}=\frac{79181}{27408}.
\]
其既约分母含素因子 \(p_\star\)（见 \eqref{eq:window6_green_kernel_audit_prime_571}），与推论 \ref{cor:xi-window6-block-tree-number-audit-prime} 的树数指纹同源。
\end{proposition}

\begin{proof}
设 \(p(t)=\prod_{j=1}^{3}(t-\lambda_j)\)（\(\lambda_j\) 为非平凡特征值）。
则 \(\frac{p'(t)}{p(t)}=\sum_{j=1}^{3}\frac{1}{t-\lambda_j}\)，取 \(t=1\) 得 \(K=\frac{p'(1)}{p(1)}\)。
将 \(p\) 的显式表达代入并化简得到 \(\frac{79181}{27408}\)（亦见证书 \eqref{eq:window6_green_kernel_audit_prime_571}）。证毕。
\end{proof}

\begin{corollary}[均值首达时间的分母谱量子化：首达时间落入 \texorpdfstring{$\mathbb{Z}[1/3,1/p_\star]$}{Z[1/3,1/p*]}]\label{cor:xi-window6-mean-first-passage-denominator-quantization}
令 \(m_{\alpha\to\beta}\) 为粗粒化链 \(T\) 上从块 \(\alpha\) 首达块 \(\beta\) 的期望步数（\(\alpha\neq\beta\)）。
则全部 \(m_{\alpha\to\beta}\) 为有理数，且其既约分母仅含素因子 \(3\) 与 \(p_\star\)；更具体地，在本 window-\(6\) 骨架中有
\[
m_{\alpha\to\beta}\in \frac{1}{3^{2}p_\star}\mathbb Z,\qquad \alpha\neq\beta,
\]
其全量闭式由表 \ref{tab:window6_mean_first_passage_times} 给出。
\end{corollary}

\begin{proof}
均值首达时间可由基本矩阵 \(Z\) 闭式表示：对 \(\alpha\neq\beta\) 有
\(
m_{\alpha\to\beta}=(Z_{\beta\beta}-Z_{\alpha\beta})/\pi_\beta
\)。
由命题 \ref{prop:xi-window6-green-kernel-arithmetic-audit-prime}，\(Z\) 的分母谱仅含 \(2,3,p_\star\)；
结合平稳分布 \(\pi_\beta=|U_\beta|/64\)（\(|U_\beta|\in\{27,22,9,6\}\)）并化简，得到 \(\alpha\neq\beta\) 时分母仅含 \(3,p_\star\)。
具体分母量子化与全量闭式由表 \ref{tab:window6_mean_first_passage_times} 可审计核验。证毕。
\end{proof}

\begin{table}[H]
\centering
\scriptsize
\setlength{\tabcolsep}{6pt}
\caption{Mean first passage times $m_{\alpha\to\beta}$ for the window-6 four-block coarse chain $T$ (exact rationals; diagonal omitted).}
\label{tab:window6_mean_first_passage_times}
\begin{tabular}{c c c c c}
\toprule
 & $U_2$ & $U_1$ & $U_L$ & $U_R$ \\
\midrule
$U_2$ & -- & $(\frac{4730}{1713})$ & $(\frac{11404}{1713})$ & $(\frac{18338}{1713})$ \\
$U_1$ & $(\frac{10834}{5139})$ & -- & $(\frac{33718}{5139})$ & $(\frac{59032}{5139})$ \\
$U_L$ & $(\frac{3764}{1713})$ & $(\frac{4718}{1713})$ & -- & $(\frac{18550}{1713})$ \\
$U_R$ & $(\frac{3310}{1713})$ & $(\frac{5768}{1713})$ & $(\frac{11162}{1713})$ & -- \\
\bottomrule
\end{tabular}
\end{table}


% AUTO-GENERATED by scripts/exp_window6_green_kernel_audit_prime_571.py
\begin{equation}\label{eq:window6_green_kernel_audit_prime_571}
\begin{aligned}
p_\star&:=571,\\
\det\bigl(I-T+\mathbf 1\pi^\top\bigr)&=\frac{9136}{8019}=\frac{2^{4} \cdot 571}{3^{6} \cdot 11},\\
K:=\sum_{\lambda\in\mathrm{Spec}(T)\setminus\{1\}}\frac{1}{1-\lambda}=\frac{79181}{27408}=\frac{79181}{2^{4} \cdot 3 \cdot 571}.
\end{aligned}
\end{equation}


\begin{proposition}[边界块 \texorpdfstring{$U_R$}{UR} 的首返时间分布完全可解：有理型母函数与显式二阶矩]\label{prop:xi-window6-UR-first-return-time-closed-form}
\leanverified{paper\_xi\_window6\_ur\_first\_return\_time\_closed\_form}
以上述四状态核 \(T\) 为动力学（\eqref{eq:window6_edge_flux_skeleton}），令 \(\tau_R\) 为从状态 \(U_R\) 出发的首返时间（首次回到 \(U_R\) 的步数，取值于 \(\mathbb N\)）。
则其概率母函数 \(F_R(z):=\mathbb E[z^{\tau_R}]\) 为有理函数，并满足闭式证书
% AUTO-GENERATED by scripts/exp_window6_ur_first_return_time.py
\begin{equation}\label{eq:window6_ur_first_return_time_pgf}
F_R(z)=\frac{z\,(110z^3+1525z^2-1221z-10692)}{623z^3+14317z^2+71010z-96228}.
\end{equation}

从而严格矩量为
\[
\mathbb E[\tau_R]=\frac{32}{3},
\qquad
\mathrm{Var}(\tau_R)=\frac{1695268}{15417}.
\]
此外，\(F_R\) 的收敛半径由分母多项式的唯一正根 \(\rho_R>1\) 决定；因此 \(\tau_R\) 具有严格指数型尾部衰减。
\end{proposition}

\begin{corollary}[首返尾率与模尺度的二重非共振：\texorpdfstring{$\log\rho_R/\log\varphi\notin\mathbb Q$}{log rhoR/logphi notin Q}]\label{cor:xi-window6-logrhoR-logphi-irrational}
在命题 \ref{prop:xi-window6-UR-first-return-time-closed-form} 的设定下，令
\[
q_R(z):=623z^3+14317z^2+71010z-96228\in\mathbb Z[z],
\]
并令 \(\rho_R>1\) 为 \(q_R\) 的唯一正根（即 \(F_R\) 的收敛半径）。
则 \(\rho_R\) 为代数次数 \(3\) 的实数，且满足
\[
\frac{\log\rho_R}{\log\varphi}\notin\mathbb Q.
\]
并且同样存在 Baker 型有效下界：存在常数 \(C_R>0\) 使对任意整数对 \((n,k)\neq(0,0)\) 有
\[
\bigl|n\log\rho_R-k\log\varphi\bigr|
\ \ge\
\exp\!\bigl(-C_R\log(1+\max\{|n|,|k|\})\bigr).
\]
\end{corollary}

\begin{proof}
由模 \(5\) 检验可知 \(q_R(z)\equiv 3z^3+2z^2+2\ (\mathrm{mod}\ 5)\) 在 \(\mathbb F_5\) 上无根，从而 \(q_R\) 在 \(\mathbb Q[z]\) 上不可约；因此 \(\rho_R\) 为三次代数数。
令 \(K_R:=\mathbb Q(\rho_R)\)，则 \([K_R:\mathbb Q]=3\)。
与推论 \ref{cor:xi-window6-field-disjointness-degree-six} 同理可得 \(K_R\cap K_\varphi=\mathbb Q\)。
若 \(\log\rho_R/\log\varphi=p/q\in\mathbb Q\)（既约），则 \(\rho_R^q=\varphi^{p}\in K_\varphi\) 且 \(\rho_R^q\in K_R\)，故 \(\rho_R^q\in K_R\cap K_\varphi=\mathbb Q\)，矛盾（因 \(\rho_R>1\Rightarrow p\neq 0\) 且 \(\varphi^{p}\notin\mathbb Q\)）。
最后一式由 Baker--W\"ustholz/Matveev 的线性对数形式下界直接推出（\cite{BakerWustholz2007LogarithmicForms,Matveev2000LinearFormsLogs}）。证毕。
\end{proof}

\begin{corollary}[三底对数的有理线性独立性与二维环面稠密性]\label{cor:xi-window6-logrhoR-logphi-logL-independent}
在推论 \ref{cor:xi-window6-logrhoR-logphi-irrational} 的设定下，对任意整数 \(L\ge 2\)，三数
\(
\log\rho_R,\ \log\varphi,\ \log L
\)
在 \(\QQ\) 上线性无关：若整数 \(m,n,k\in\ZZ\) 满足
\[
m\log\rho_R+n\log\varphi+k\log L=0,
\]
则必有 \(m=n=k=0\)。
\par
因此 \(\log\rho_R/\log L\notin\QQ\)，从而一参数轨道
\[
\bigl\{\,t(\log\rho_R,\log L)\bmod 2\pi:\ t\in\RR\,\bigr\}
\]
在二维环面 \((\RR/2\pi\ZZ)^2\) 上稠密；等价地，任何同时依赖两底相位 \(t\log\rho_R\) 与 \(t\log L\) 的有限型谱构造都不可能出现非零的共同纯虚周期。
\end{corollary}
\begin{proof}
若 \(m\log\rho_R+n\log\varphi+k\log L=0\)，指数化得
\(
\rho_R^{m}\varphi^{n}L^{k}=1
\)，即
\(
\rho_R^{m}=\varphi^{-n}L^{-k}\in\QQ(\varphi)
\)。
左端属于 \(\QQ(\rho_R)\)，故 \(\rho_R^{m}\in \QQ(\rho_R)\cap\QQ(\varphi)=\QQ\)（推论 \ref{cor:xi-window6-logrhoR-logphi-irrational} 的证明已给出交域为 \(\QQ\)）。
由于 \(\rho_R\) 为三次代数数且 \(\rho_R\notin\QQ\)，只能有 \(m=0\)。
此时 \(\varphi^{n}L^{k}=1\)。又 \(L^{k}\in\QQ\)，故 \(\varphi^{n}\in\QQ\)，从而 \(n=0\)，继而 \(k=0\)。
稠密性为 Kronecker--Weyl 定理在二维情形的直接推论：\(\log\rho_R/\log L\notin\QQ\) 等价于二维旋转稠密。证毕。
\end{proof}

\begin{corollary}[有效有理逼近下界与幂律锁定排除]\label{cor:xi-window6-alphaR-rational-approx-bound}
在推论 \ref{cor:xi-window6-logrhoR-logphi-irrational} 的设定下，令
\[
\alpha_R:=\frac{\log\rho_R}{\log\varphi}.
\]
则对任意整数 \(n\ge 1\) 与任意 \(k\in\ZZ\)，成立显式的有理逼近下界
\[
\Bigl\lvert \alpha_R-\frac{k}{n}\Bigr\rvert
=\frac{\bigl|n\log\rho_R-k\log\varphi\bigr|}{n\log\varphi}
\;\ge\;
\frac{1}{n\log\varphi}\exp\!\Bigl(-C_R\log\bigl(1+\max\{n,|k|\}\bigr)\Bigr).
\]
从而不存在非零整数对 \((n,k)\) 使 \(\rho_R^n=\varphi^k\)；
并且任何试图在黄金尺度 \(\varphi^k\) 与首返尾率尺度 \(\rho_R^n\) 之间建立精确幂律锁定（或误差可忽略的近幂律锁定）的方案，均受该有效非共振下界约束。
\end{corollary}
\begin{proof}
由推论 \ref{cor:xi-window6-logrhoR-logphi-irrational} 的 Baker 型下界，直接除以 \(n\log\varphi\) 得到所示不等式。
若存在 \((n,k)\neq(0,0)\) 使 \(\rho_R^n=\varphi^k\)，则 \(n\log\rho_R-k\log\varphi=0\)，与下界矛盾，故不存在幂律锁定。证毕。
\end{proof}

\begin{corollary}[\texorpdfstring{$(1+3+8)$}{(1+3+8)} 的倒数通量：边界稳态通量常数恰为 \texorpdfstring{$1/\dim\mathfrak g_{\mathrm{SM}}$}{1/dim gSM}]\label{cor:xi-window6-reciprocal-flux-boson12}
\leanverified{paper\_xi\_window6\_reciprocal\_flux\_boson12}
在 window-\(6\) 的四块骨架证书中（\eqref{eq:window6_edge_flux_skeleton}），边界块 \(U_R\) 的割大小与稳态通量常数满足
\[
\Phi_R=\frac{|\partial U_R|}{64\cdot 6}=\frac{1}{12}.
\]
若将标准模型规范玻色子计数解释为 \(\dim\mathfrak g_{\mathrm{SM}}=1+3+8=12\)，则得到严格倒数律
\[
\Phi_R=\frac{1}{\dim\mathfrak g_{\mathrm{SM}}}.
\]
\end{corollary}

\begin{corollary}[复位事件时钟的分支频率刚性：\texorpdfstring{$\{34,55\}$}{\{34,55\}} 的极限频率由 \texorpdfstring{$\varphi^{-1}$}{phi^{-1}} 唯一决定]\label{cor:xi-window6-reset-clock-gap-frequencies}
\leanverified{paper\_xi\_window6\_reset\_clock\_gap\_frequencies}
对 window-\(6\) 的复位事件集合 \(\mathcal R_6\)（定义 \ref{def:terminal-reset-events}），其返回时间二值谱为
\(\{F_9,F_{10}\}=\{34,55\}\)（定理 \ref{thm:terminal-reset-events-sturmian}；并见表 \ref{tab:fold_zeck_reset_event_gaps_m6}）。
记字母 \(A:=55\)、\(B:=34\)，则差分字母序列为 Fibonacci 替换 \(\sigma(A)=AB,\sigma(B)=A\) 的不动点；因此两间隔的极限频率存在且为
\[
\lim_{N\to\infty}\frac{\#\{n\le N:\ r_{n+1}-r_n=55\}}{N}=\frac{1}{\varphi},
\qquad
\lim_{N\to\infty}\frac{\#\{n\le N:\ r_{n+1}-r_n=34\}}{N}=\frac{1}{\varphi^2}.
\]
从而 “\(SO(10)\) 尾部偏移 \(34\)” 在无限时钟尺度上获得不可调的概率权重，为统一分支的概率化选择给出可审计的确定性频率定律。
\end{corollary}

\begin{theorem}[尾部偏移半群的 Frobenius 阈值：\texorpdfstring{$\langle 21,34,55\rangle$}{<21,34,55>} 的导通界]\label{thm:xi-window6-tail-semigroup-frobenius-threshold}
\leanverified{paper\_xi\_window6\_tail\_semigroup\_frobenius\_threshold}
在 window-\(6\) uplift 的离散尾部偏移口径下，非零偏移仅可能取自
\(\{21,34,55\}=\{F_8,F_9,F_{10}\}\)（见 \S\ref{subsubsec:window6-terminal-propositions} 的相应结论链）。
令数值半群
\[
S_6:=\langle 21,34,55\rangle\subset\mathbb N.
\]
则 \(S_6\) 为余有限半群，且其 Frobenius 数与导通数满足
% AUTO-GENERATED by scripts/exp_window6_tail_semigroup_frobenius.py
\begin{equation}\label{eq:window6_tail_semigroup_frobenius}
g(S_6)=659,\qquad c(S_6)=660.
\end{equation}

换言之，对任意整数 \(n\ge 660\) 都存在 \(a,b,c\in\mathbb N\) 使
\(
n=21a+34b+55c
\)，而 \(659\) 为最大不可表示数。
\end{theorem}

\begin{corollary}[维数筛法的有限性：在尾部可加假设下 \texorpdfstring{$D>671$}{D>671} 不能再由维数否证]\label{cor:xi-window6-dimension-sieve-finite}
\leanverified{paper\_xi\_window6\_dimension\_sieve\_finite}
若在“额外规范模态仅通过尾部偏移的可加叠加进入”这一可审计假设下，将可达玻色子总数集合写为
\[
\mathcal D_{\mathrm{reach}}:=12+S_6,
\]
则由定理 \ref{thm:xi-window6-tail-semigroup-frobenius-threshold} 的导通数立得：
所有 \(D\ge 12+660=672\) 均属于 \(\mathcal D_{\mathrm{reach}}\)，而不可达维数集合 \(\mathbb N\setminus\mathcal D_{\mathrm{reach}}\) 为有限集，且其最大元为 \(12+659=671\)。
因此，超过 \(671\) 的维数不能再仅由尾部可达性被否证；必须引入更细的签名 \(\Sigma\)、退化谱、或谱穿越证书等更高分辨率不变量。
\end{corollary}

\begin{proposition}[从半群阈值到可证伪预测：硬否证与软筛选的分层]\label{prop:xi-window6-falsifiability-layering-from-semigroup}
\leanverified{paper\_xi\_window6\_falsifiability\_layering\_from\_semigroup}
在区间 \(D\le 671\) 内，\(\mathcal D_{\mathrm{reach}}\) 之外的维数点构成绝对否证点：任何声称由 window-\(6\) 尾部可加机制产生的统一候选若具有这些维数，则与协议的可达性约束直接矛盾。
相反，对 \(\mathcal D_{\mathrm{reach}}\cap[0,671]\) 内的维数点，维数本身不再构成否证，必须进一步满足 Zeckendorf--Fibonacci 分辨率签名的一致审计（例如禁相邻约束、边界层激活的无冗余性以及谱/PSD 证书的兼容性）。
因此可证伪性在该框架中被精确分割为：低维有限区间内由半群阈值提供的硬否证，与同区间内由签名/谱证书提供的软筛选两类。
\end{proposition}
